\epigraph{The brain is a wonderful organ; it starts working the moment you get up in the morning and does not stop until you get into the office}
{\textit{Robert Frost}}

\section{Overview}

% Talk about how lectures are gonna be different to classroom learning and how to deal with that, how the material comes up quite fast. 
% how many lectures per module 
% what is the general structure like of lectures 
% Note on notes being provided beforehand + other resources 
% how long they are 

In your first year, your typical week will consist of a mixture of lectures, seminars, tutorials, problem classes, labs, coursework and workshops. 

For example, in my first semester, I had four modules consisting of three lectures a week and a problem class lasting 45 minutes. Alongside this, I had Core Skills and workshop sessions, which left me plenty of time to be involved in a range of clubs and societies.

\subsection{Core Skills}
% info on core skills module? - what it's about..? - basically mention infor that is alreaady on blackboard 
Core Skills is a compulsory and non-credited module that is taken by everyone in the first year. 

As part of the module, you will have to pass two Core Maths Skills Tests, one per semester. The tests gives you a way of making sure you have these skills right from the beginning of your study (Semester 1) and as your study progresses (Semester 2). Both of these are online exams that can be taken repeatedly until you have passed. 

There are also scheduled sessions on employability and digital skills. These sessions are designed to maximise your digital literacy and employment prospects during and after your degree. 
% links to careers page, etc. 


\subsection{Problem Classes}
% What they are - how to make the most of them 
% Try and have a look at the problems beforehand 
% what is the mode up-to-date info on this? 

Each module will have a weekly problem class. These are quite different from lectures. Instead of listening to new material being introduced, the focus is on working through questions and understanding how the ideas from the lectures actually get used.

The exact format can vary slightly between modules and lecturers. Sometimes the lecturer will go through selected problems on the board, sometimes students will suggest questions they want to see explained, and sometimes you go might go through some coursework questions. 

One of the most important thing you can do to make problem classes useful is look at the problems beforehand. If you walk into a problem class having never seen the sheet before, it’s much harder to follow what’s going on and becomes easy to just copy down solutions without really understanding them.

\section{Online Learning Tools}

% timetable can be found on Seats and on blackboard - useful to link timetable to google calendar or something equivalent. 
% Check in with Seats - QR code or the seats code will usually be given start of the lecture  - keeps a good track of attendance 


You'll become familiar with a variety of names very quickly: \textbf{Blackboard}, \textbf{Seats}, and \textbf{Panopto}

Blackboard is the university's online learning platform where lectures notes are uploaded, announcements are posted and the coursework and recordings are stored. If something official happens in the module, it will most likely appear there. It's useful having a quick look at notes before the lecture if you can, even just seeing the headings can make it much easier to follow what’s happening when you’re sitting in the lecture hall.

Panopto is where you can access lecture recordings, including historical recordings. 

One of the most useful things you can do in the first week is link your timetable to your own calender, whether that's Google Calendar or something equivalent. Remember that universities will not chase you the way school did. There are no bells and no-one reminds you every morning where you're supposed to me, so the responsibility shifts to you. 

Seats is the system used for timetabling and attendance. It shows your lecture locations, times, and room numbers. It's also how attendance is recorded. At the start of lectures, a QR or Seats code will be displayed on the board which students can scan or type to register their attendance on the app. 

\section{Recorded Doesn't Mean Optional}

You will quickly realise that most lectures are recorded, and at some point you might be temped to think, `I can just watch it later', especially on a cold and rainy morning, or after a late night, or when the idea of sitting in a lecture feels mildly exhausting. 

I cannot stress this enough, but being physically present in a lecture matters more than you think. When you are in the room, you are less likely to pause a video and get distracted by your phone. You are also less likely to tell yourself that you'll `finish it later' and risk not coming back to it. You are more likely to stay with the explanations as they unfold, even if you don't fully understand them yet. You hear the throwaway comments that never make it into typed notes, and pick up on the emphasis in the lecturer's voice when they say something is important, and that emphasis stays with you. 

There is also something about the collective focus that recordings cannot replicate. Sitting in a room full of people who are all attempting to follow the same proof creates a concentration that watching it alone at 1.5x speed, half distracted cannot compete with. In addition to that, you are able to ask questions during or after the lecture if you attend in person. 

Recordings are useful, in fact, they are incredibly useful for revision or revisiting a proof if that felt unclear the first time. They are a safety net, but not a replacement for showing up. 

\section{Catching Up On Lectures}

At some point, you may very well miss a lecture. Maybe you were ill, overslept, or just needed a morning off. It happens, and one missed lecture is not a crisis.

The trick to avoid spiralling is to catch up quickly. If you miss a lecture, watch the recording within a day or two while the topic and notation is still familiar. Don't let it all pile onto future you.  

% SHORT ASIDE ON FUNNY WAY OF GETTING YOURSELF BACK INTO FOCUS, LIKE HOW I HAD TO BE IN THE GYM DOING A 2 hr ZONE 2 RIDE WHILE WATCHING CALC RECORDINGS. 

When you do catch up, don't treat the recording like a background noise. Try to follow the argument yourself. Stop and ask whether you understand what each part is doing. 

You won't always find it easy to sit down and follow a lecture, I certainly couldn't at all times. In fact, at one point in Semester 2, I'd go to the gym on weekends, get on one of the stationary bikes and rewatch my Calculus lectures while cycling. I wasn't taking notes or anything, but I was surprisingly engaged with the material. There isn't a single `correct' way to study and sometimes, the hardest part is just getting yourself to engage with the material at tall, so if something slightly unconventional helps you do that, it's probably worth trying. 

\section{Asking questions}
% Why it's a good thing - link back to being present in person 
% Addressing fear and embarassemtn of asking questions in lectures 

If you ever find yourself sitting in a lecture room, slightly unsure about a definition or confused about some step in a proof, and you look around the room and everyone else looks calm, and you might think, `Maybe it's just me'. 

It almost never is. Most people are running through the same internal debate of `Is this worth asking? Is this obvious? Is this too stupid of a question?', etc. 

Questions like `How did we go from Line 1 to Line 2 of this proof', or `Is there an example where this conditions fails?' are perfectly good questions to ask in lecture. There are no stupid questions. 

By asking questions, you benefit yourself and others, including the lecturer. Learning mathematics is a process and asking questions is part of that process. If asking during the lecture feels too intimidating, most lecturer's are more than happy for students to ask them at the end of the lecture, or during office hours and workshops. 

\subsection{Talking in lectures}
% Basically - pls shut the fuck up and dont' talk over lectueres, it's rude, distruptive, interrupts other people's focus 

Lectures are shared learning spaces and everyone in the room is trying to concentrate on the same material. It might seem harmless to chat quietly during a lecture, but even small conversations can be a distracting for the people sitting nearby especially with an attention demanding subject like maths. 

More importantly, it's also a matter of respect. The lecturer has prepared the material and everyone else in the room has chosen to be there to learn. Talking over them or holding side conversations can come across as rude even if that isn't the intention. 

If you want to discuss something with the person next to you, it's usually better to wait for a natural pause, write it down somewhere, or bring it up after the lecture. When everyone treats the lecture space with a bit of consideration, it makes it much easier for everyone else to follow what's going on. 


\section{Giving feedback}

% Have more people do feedback 
% suggsetion box in the studetn centre - 
% talk to your course reps 

At some point during the semester, you will be asked to fill in feedback forms for SSLC or mid-semester. 

It is very tempting to either rush through them, or completely ignore them. It is also very easy to write things like `Too fast' or `Confusing' or `Good'. 

Feedback only becomes useful when it is specific. Writing `This lecture is confusing' doesn't give anyone anything to work with, whereas writing `There are too many definitions and very few examples to see how those definitions work' is useful. 

Some responses convey reaction, others information, so ensure that whenever you're in a position of providing feedback, you relay the information over reaction as best as you can. 

Lecturers are not mind readers. They often assume a certain level of clarity because they've been living with the material for far far longer than we have. So if ten students feel quietly lost but say nothing, nothing changes. 

Feedback not only helps the lecturers, but also the cohort. And it helps the individual, because the act of explaining why something was unclear forces you to reflect on what you were expecting. 

\subsection{Too fast? Too Slow?}

TAt some point in the semester, you’ll be asked to complete SSLC or mid-semester feedback forms. It's very tempting to either rush through them and give vague feedback, or completely ignore them. 

Feedback only works when many students contribute and when what they write is specific. One vague comment does nothing, but ten clear comments pointing to the same issue are hard to ignore. 

Bad feedback describes a reaction, good feedback gives information someone can act on. Writing `confusing lecture' doesn't give anyone anything to work with, whereas writing `There are too many definitions and very few examples to see how those definitions work' or `slides move on before there's time to copy or process examples is useful. 

Lecturers are not mind readers. They’ve worked with the material for years, so what feels obvious to them may not be obvious to you. If many students are quietly lost but say nothing, nothing changes. Feedback helps the lecturers who can adjust delivery or structure which turn benefits the cohort, and it also helps you because explaining confusion forces you to think about what you didn’t understand. 

\subsection{Pace of lectures}

Some lectures will feel too fast. This often happens when you’re seeing abstract ideas for the first time. That discomfort is normal as understanding usually comes after repeated exposure and not always during the first pass.

Other lectures will feel slow, perhaps even repetitive, and you may find yourself zoning out because you think you already understand it. Instead of disengaging, use that time to look deeper into what it being taught, asking why the definition is phrased in a certain way, or what breaks if an assumption is removed, or even how the idea connects to something you've already seen. 

The pace of a lecture may not match your personal learning speed exactly, and that's normal. But if something feels consistently unreasonable, say so clearly and specifically in feedback or directly to the lecturer.


\section{Coursework}
% Make a note of when it is due - most modules follow a timetable, eg: setting work every friday and due the next friday 

Most modules include coursework, usually in the form of weekly problem sheets. These are typically released on a fixed day and due the following week, although the exact timing varies by module. Some are submitted on paper, others through online systems.

At first, coursework can feel relentless. You submit one assignment and the next appears almost immediately, which can make it seem like there is never a clear break. In practice, this regular rhythm is intentional as it keeps you engaged with the material and prevents topics from being pushed aside and forgotten.

The key is to treat coursework as part of your weekly structure rather than as something extra. Plan your week around it. Decide early when you will start each sheet and when you aim to make progress on it, instead of leaving it to fit around everything else.

The hardest part is simply starting. A common pattern is to ignore the sheet for a few days, glance at it midweek, and then realise late in the weekend that nothing has been attempted. By that point, you are working against time rather than thinking properly about the problems. Starting earlier changes this completely. Coursework is designed to take time because it is meant to make you think, not just produce answers. Leaving it until the night before removes the opportunity to get stuck, reflect, and come back with a clearer idea.

Aim to begin at least a day or two after it is released, even if you only attempt the first question. If you get stuck after making a genuine attempt, that is exactly what the support structures are there for. 



